\documentclass[11pt,a4paper]{article}
\usepackage[utf8]{inputenc}
\usepackage[T1]{fontenc}
\usepackage{amsfonts}
\usepackage[french]{babel}
\frenchbsetup{StandardLists=true} % à inclure si on utilise \usepackage[francais]{babel}
\usepackage{enumitem}
\usepackage{amssymb}
\usepackage{amsmath}
\usepackage[left=2cm,right=2cm,top=2cm,bottom=2cm]{geometry}
\usepackage{multirow}
\usepackage[dvipsnames]{xcolor}

\usepackage[T1]{fontenc}
\usepackage{fouriernc}
\usepackage{graphicx}

\usepackage{setspace}
\singlespacing

\usepackage{pstricks}
\usepackage{xkeyval}
\usepackage{pst-labo}
\usepackage{chemist} %chimie


\usepackage{multicol}
\usepackage{caption}
\usepackage{fancyhdr}
\pagestyle{fancy}
\renewcommand\headrulewidth{1pt}
\fancyhead[L]{Lycée Denis Diderot }
\fancyhead[R]{Classe de 1 Bac Pro }
\setcounter{page}{2}\fancyfoot[C]{page \thepage}
\fancyfoot[R]{Année 2024-2025}
\fancyfoot[L]{Alexis RUIZ}
\usepackage{multirow,tabularx,booktabs}
\usepackage[tikz]{bclogo} 

\newcommand{\cadreblanc}[1]{%
\medskip\framebox[\linewidth][c]{%
\rule{0mm}{#1mm}}\par
\medskip}

\usepackage{awesomebox}
\setlength{\aweboxvskip}{0mm}
\usepackage{pifont}
\usepackage{chemmacros}
\chemsetup{modules=all}
\setlength{\columnseprule}{.25mm}

\renewcommand{\arraystretch}{2.5}
\newcommand*\acocher{\quitvmode{\fboxsep0pt \fboxrule0.8pt \fbox{\vrule height1.8ex depth.1ex width0pt \vrule height0pt depth0pt width1.9ex }}}

\newcolumntype{C}[1]{>{\centering\arraybackslash }b{#1}}

\newcommand{\code}{\awesomebox[black]{2pt}{

\includegraphics[scale=0.15]{Chapitre 1 - Energie et puissance électrique/images/frame.png}
}{black}}

\newcommand{\moteur}{\awesomebox[black]{2pt}{

\includegraphics[scale=1]{images/moteur.jpg} 
}{black}}

\newcommand{\defconclusion}{\awesomebox[black]{2pt}{\rotatebox{90}{\Large{\textbf{Conclusion}}}}{black}}

\newcommand{\defbox}{\awesomebox[black]{2pt}{
\includegraphics[scale=0.5]{Chapitre 1 - Energie et puissance électrique/images/appel exam.png}}{black}}

  \newenvironment{itemize*}%  % listes compactes
    {\begin{itemize}%
      \setlength{\itemsep}{0pt}%
      \setlength{\parskip}{0pt}}%
    {\end{itemize}}

\frenchbsetup{StandardLists=true}


\newcommand{\cadre}[1]
{\begin{bclogo}[couleur = white,logo=\bcinfo , arrondi = 0.1]
{
~~
{#1}
~~
\end{bclogo}}}



\begin{document}

\flushleft \textbf{NOM : \hfill Classe : \hspace{3cm} ~~\\[2mm]
Prénom:\\[0,4cm]
}
\hrule\
\\[0,2cm]

\begin{center}
 \LARGE{Activité - Dissolution et dilution   }
\end{center}

\vspace{3 mm}

\normalsize

\hrule\
\\[0,2cm]
   \begin{center}
   \textbf{Le soin et la rédaction interviendront dans la notation.}
   \end{center} 





\begin{bclogo}[couleur=white,  arrondi =0.1 , logo = \bcinfo, epBarre =2]{ Objectifs }{
\begin{itemize*}
\item Réaliser une solution de concentration donnée par dilution ou par dissolution;
\item Calculer une masse molaire moléculaire à partir du tableau périodique et de la formule de la molécule;
\item Calculer la concentration d'un soluté.
\end{itemize*}
}
\end{bclogo}

\vfill 


\fbox{\large{\ding{202} Dissolution et dilution : }}

\vfill 

\begin{bclogo}[couleur=white,  arrondi =0.1 , logo = \bcinfo, epBarre =2]{A retenir }
\vspace{0.1cm}

\begin{itemize}
\item Une \textbf{dissolution} consiste à préparer une solution à partir d'un soluté solide.

\begin{center}
\includegraphics[scale=0.78]{Chapitre 4 - Solutions aqueuses/images/Disolution.png}
\vspace{0.1cm}
\end{center}
\end{itemize}
\end{bclogo}


\vfill 

\newpage

\begin{bclogo}[couleur=white,  arrondi =0.1 , logo = \bcinfo, epBarre =2]{A retenir }
\vspace{0.1cm}

\begin{itemize}
\item Une \textbf{dilution} consiste à préparer une solution à partir d'une autre solution (en diminuant sa concentration)

\begin{center}
\includegraphics[scale=0.65]{Chapitre 4 - Solutions aqueuses/images/Dilution2.png}
\vspace{0.1cm}
\end{center}
\end{itemize}
\end{bclogo}





\textbf{Compléter} la liste du matériel nécessaire pour chaque manipulation.  \\

\begin{center}
\begin{tabular}{|c|c|}
\hline
~~~~~~~~~~~~~~~~~~~~~~~~~~~~Dissolution~~~~~~~~~~~~~~~~~~~~~~~~~~~~ & ~~~~~~~~~~~~~~~~~~~~~~~~~~~~Dilution ~~~~~~~~~~~~~~~~~~~~~~~~~~~~\\
\hline
&  \\
\hline
&  \\
\hline
& \\
\hline
& \\
\hline
& \\
\hline
\end{tabular}
\end{center}


\fbox{\large{\ding{203} Dissolution et dilution : }}

Suite à une consultation chez le médecin pour de nombreuses et répétitives crampes, un adulte obtient les résultats partiels d'une analyse de sang faite pour une personne dont on surveille la glycémie : 
\begin{center}
\fbox{\includegraphics[scale=0.7]{Chapitre 4 - Solutions aqueuses/images/gly.png}}
\end{center}

\newpage

\begin{bclogo}[couleur=white,  arrondi =0.1 , logo = \bcinfo, epBarre =2]{Problématique :  }
Comment déterminer si la glycémie de cette personne est dans la normale ou non ?
\end{bclogo}


\begin{enumerate}
\item  \textbf{Donner} la valeur et l'unité de la glycémie de la personne :

\vspace{3mm}
\dotfill 

\item  \textbf{Donner} les valeurs de référence de glycémie ainsi que leur unité : 

\vspace{3mm}
\dotfill 

\item  \textbf{Expliquer} pourquoi on ne peut pas répondre à la problématique à ce stade :

\vspace{3mm}
\dotfill 

\vspace{3mm}
\dotfill 

\end{enumerate}

\vspace{0.1cm}

\begin{bclogo}[couleur=white,  arrondi =0.1 , logo = \bcinfo, epBarre =2]{ Objectifs }{
La glycémie contrôle le taux de glucose de formule chimique \chemform{C_6H_{12}O_6} dans le sang d'une personne. Les chimistes expriment la quantité de matière par "paquets" d'atomes ou de molécules que l'on appelle "moles" (le symbole de cette unité est \textbf{mol}) 
}
\end{bclogo}

Pour trouver la masse d'une mole de molécule de glucose (un paquet) \textbf{on additionne la masse molaire de tous les atomes}. Cette masse molaire se trouve dans le tableau périodique :
\begin{center}
\fbox{\includegraphics[scale=0.4]{Chapitre 4 - Solutions aqueuses/images/massemol.png}}
\end{center}

On donne ci-dessous un extrait de classification périodique avec les composants du glucose :
\begin{center}
\begin{tikzpicture}
\draw (0,0)--(0,3)--(2,3)--(2,0)--cycle;
\node at (1,1.75) {{\LARGE H}};
\node at (1,0.3) {1,0};
\node at (1,1) {Hydrogène};
\node at (0.3,2.6){{\Large 1}};
\end{tikzpicture}
\hspace{1cm}
\begin{tikzpicture}
\draw (0,0)--(0,3)--(2,3)--(2,0)--cycle;
\node at (1,1.75) {{\LARGE C}};
\node at (1,0.3) {12,0};
\node at (1,1) {Carbone};
\node at (0.3,2.6){{\Large 6}};
\end{tikzpicture}
\hspace{1cm}
\begin{tikzpicture}
\draw (0,0)--(0,3)--(2,3)--(2,0)--cycle;
\node at (1,1.75) {{\LARGE O}};
\node at (1,0.3) {16,0};
\node at (1,1) {Oxygène};
\node at (0.3,2.6){{\Large 8}};
\end{tikzpicture}
\end{center}

\begin{enumerate}
\item \textbf{Calculer},  en g/mol, la masse molaire moléculaire du glucose (\chemform{C_6H_{12}O_6}) :

\vspace{3mm}
\dotfill 

\item \textbf{En utilisant une formule présente dans votre cours}, répondre à la problématique (attention aux unités, il faudra transformer avant les \si{\milli\mol} en \si{mol})  : 

\vspace{3mm}
\cadreblanc{20}

\end{enumerate}

\vspace{1cm}

\newpage


\begin{center}
 \LARGE{FICHE METHODE  }
\end{center}

\vspace{3 mm}

\normalsize

\hrule\


\begin{bclogo}[couleur=white,  arrondi =0.1 , logo = \bcinfo, epBarre =2]{A retenir }
\vspace{0.1cm}

\begin{itemize}
\item Une \textbf{dissolution} consiste à préparer une solution à partir d'un soluté solide.

\begin{center}
\includegraphics[scale=0.78]{Chapitre 4 - Solutions aqueuses/images/Disolution.png}
\vspace{0.1cm}
\end{center}
\end{itemize}
\end{bclogo}

\vfill 

\begin{bclogo}[couleur=white,  arrondi =0.1 , logo = \bcinfo, epBarre =2]{A retenir }
\vspace{0.1cm}

\begin{itemize}
\item Une \textbf{dilution} consiste à préparer une solution à partir d'une autre solution (en diminuant sa concentration)

\begin{center}
\includegraphics[scale=0.8]{Chapitre 4 - Solutions aqueuses/images/Dilution2.png}
\vspace{0.1cm}
\end{center}
\end{itemize}
\end{bclogo}


\vfill 




\end{document}
